\documentclass{article}
	\usepackage{graphicx}%插入图片
	\usepackage{titlesec}%修改标题格式
	\usepackage{amsmath}%大部分数学符号
	\usepackage{mathrsfs}%花体
	\usepackage{amssymb}%双线体mathbb
	\usepackage{xeCJK}%中文支持
	\usepackage{ctex} %中文字体
	\usepackage{geometry}%几何绘图
	\usepackage{setspace}%设置行距
	\usepackage{multicol}%表格合并
	\usepackage{multirow}%表格合并
	\usepackage{subfigure}%次级图片
	\usepackage{hyperref}%调整标注与引用
	\usepackage{indentfirst}%首行缩进}
\usepackage{tikz}%tikz绘图
\usepackage{tikz-feynman}
\usepackage{listings}%代码插入
\lstset{language=Matlab}
\usepackage{xcolor}
\usepackage{longtable}%大型表格跨页显示
\usepackage{booktabs}
% \usepackage{mathpazo} %palatino equations
\usepackage{pgfplots}
\usepackage{siunitx}
\usepackage{mhchem}
\usepackage{slashed}
\linespread{1.25}
\begin{document}
\begin{flushright}
    规范场
    
    hw2

    2024310867 王亚朋
\end{flushright}
\begin{center}
    \Large{\textbf{Homework 3}}
\end{center}

\section*{Problem 1: One-loop RG for complex scalar}

\begin{equation}
\mathcal{L} = \partial_\mu\phi\,\partial^\mu\phi^* - m^2|\phi|^2 - \frac{\lambda}{2}(|\phi|^2)^2.
\end{equation}
Decompose $\phi = (\phi_1 + i\phi_2)/\sqrt{2}$:
\begin{equation}
\mathcal{L} = \frac12(\partial\phi_1)^2 + \frac12(\partial\phi_2)^2 - \frac{m^2}{2}(\phi_1^2+\phi_2^2) - \frac{\lambda}{8}(\phi_1^2+\phi_2^2)^2.
\end{equation}
Bare Lagrangian in $d=4-\epsilon$:
\begin{equation}
\mathcal{L}_B = Z_\phi|\partial\phi|^2 - Z_\phi Z_m m^2|\phi|^2 - \frac{\lambda}{2}\mu^\epsilon Z_\lambda Z_\phi^2(|\phi|^2)^2,
\end{equation}
$Z_i=1+\delta_i$. Counterterms:
\begin{equation}
\mathcal{L}_{\rm ct} = \delta_\phi|\partial\phi|^2 - (\delta_\phi+\delta_m)m^2|\phi|^2 - \frac{\lambda}{2}\mu^\epsilon(\delta_\lambda+2\delta_\phi)(|\phi|^2)^2.
\end{equation}

\medskip\noindent\textbf{Self-energy} (tadpole $\propto \int d^dk/(k^2-m^2)=0$ in dim reg):
\begin{equation}
-i\Sigma(p^2) = \lambda\!\int\!\frac{d^dk}{(2\pi)^d}\frac{i}{k^2-m^2} = \frac{\lambda m^2}{16\pi^2}\Bigl(\frac{2}{\bar\epsilon}+1-\ln\frac{m^2}{\mu^2}\Bigr),\quad \bar\epsilon^{-1}\equiv\epsilon^{-1}-\gamma_E+\ln4\pi.
\end{equation}
Hence
\begin{equation}
\delta_\phi = 0,\qquad \delta_m = \frac{\lambda}{16\pi^2}\frac{2}{\bar\epsilon},\qquad \gamma_m = -\frac12\frac{d\ln Z_m}{d\ln\mu} = \frac{\lambda}{16\pi^2},\qquad \gamma_\phi=0.
\end{equation}

\noindent\textbf{Four-point function.} $s$-, $t$-, $u$-channel bubbles (each symmetry factor $1/2$):
\begin{equation}
i\mathcal{M}_{1\text{-loop}} = -i\lambda - i\frac{\lambda^2}{32\pi^2}\bigl[I(s)+I(t)+I(u)\bigr],\quad I(P^2)=\frac{2}{\bar\epsilon}+\text{finite}.
\end{equation}
At $s=t=u=\mu^2$:
\begin{equation}
\delta_\lambda = \frac{5\lambda}{16\pi^2}\frac{1}{\bar\epsilon}.
\end{equation}
From $\lambda_0=\mu^\epsilon Z_\lambda\lambda$:
\begin{equation}
\beta(\lambda) = \mu\frac{d\lambda}{d\mu} = -\epsilon\lambda - \lambda\mu\frac{d\ln Z_\lambda}{d\ln\mu}
= \frac{5\lambda^2}{16\pi^2}+\mathcal{O}(\lambda^3).
\end{equation}
\begin{equation}
\boxed{\beta(\lambda)=\frac{5\lambda^2}{16\pi^2},\qquad \gamma_m=\frac{\lambda}{16\pi^2},\qquad \gamma_\phi=0.}
\end{equation}

\section*{Problem 2: QED counterterms and BPHZ}

\begin{equation}
\mathcal{L} = -\frac14 F_{\mu\nu}F^{\mu\nu} + \bar\psi(i\slashed\partial-m)\psi + e\bar\psi\slashed A\psi.
\end{equation}
Bare:
\begin{equation}
\mathcal{L}_B = -\frac14 Z_3 F^2 + Z_2\bar\psi i\slashed\partial\psi - Z_0 m\bar\psi\psi + e\mu^{\epsilon/2}Z_1\bar\psi\slashed A\psi.
\end{equation}

One-loop counterterms ($\overline{\rm MS}$):
\begin{align}
\delta_3 = Z_3-1 &= -\frac{e^2}{12\pi^2}\frac{1}{\bar\epsilon},\\
\delta_2 = Z_2-1 &= -\frac{e^2}{16\pi^2}\frac{1}{\bar\epsilon},\\
\delta_0 = Z_0-1 &= -\frac{3e^2}{8\pi^2}\frac{1}{\bar\epsilon}.
\end{align}
Ward identity $\Rightarrow Z_1=Z_2$.

RG equations:
\begin{align}
\beta(e)   &= \frac{e^3}{12\pi^2},\\
\gamma_m   &= \frac{3e^2}{8\pi^2},\\
\gamma_\psi &= \frac{e^2}{16\pi^2},\\
\gamma_A   &= -\frac{e^2}{12\pi^2}.
\end{align}

\section*{Problem 3: SU($N$) Yang--Mills $\beta$ function}

\begin{equation}
\mathcal{L}_{\rm YM} = -\frac14 F_{\mu\nu}^a F^{a\mu\nu} + \sum_{f=1}^{N_F}\bar\psi_f(i\slashed D-m_f)\psi_f.
\end{equation}
$g_0=\mu^{\epsilon/2}Z_g g$. One-loop renormalization constants:
\begin{align}
Z_3 &= 1 - \frac{g^2}{16\pi^2}\Bigl(\frac53 C_A - \frac43 T_F N_F\Bigr)\frac1{\bar\epsilon},\\
Z_1 &= 1 - \frac{g^2}{16\pi^2}\Bigl(\frac23 C_A + \frac43 T_F N_F\Bigr)\frac1{\bar\epsilon},
\end{align}
$C_A=N$, $T_F=1/2$. $Z_g=Z_1 Z_3^{-3/2}$:
\begin{equation}
Z_g = 1 + \frac{g^2}{16\pi^2}\Bigl(-\frac{11}{6}N + \frac13 N_F\Bigr)\frac1{\bar\epsilon}.
\end{equation}
\begin{equation}
\boxed{\beta(g) = -\frac{g^3}{16\pi^2}\Bigl(\frac{11}{3}N - \frac23 N_F\Bigr)}.
\end{equation}

\section*{Problem 4: Goldstone scattering equivalence}

\begin{equation}
\mathcal{L} = \partial_\mu\phi\,\partial^\mu\phi^* + m^2|\phi|^2 - \frac{\lambda}{2}(|\phi|^2)^2,\qquad v = \frac{m}{\sqrt{\lambda}}.
\end{equation}

\noindent\textbf{Linear:} $\phi = v + (\phi_1+i\phi_2)/\sqrt{2}$.
\begin{equation}
|\phi|^2 = v^2 + \sqrt{2}v\phi_1 + \frac{\phi_1^2+\phi_2^2}{2},\qquad
-\frac{\lambda}{2}(|\phi|^2)^2 \supset -\frac{\lambda}{8}\phi_2^4 = -\frac{3\lambda}{4!}\phi_2^4.
\end{equation}
\begin{equation}
i\mathcal{M}_{\phi_2\phi_2\to\phi_2\phi_2} = -i(3\lambda).
\end{equation}

\noindent\textbf{Polar:} $\phi = (v+R/\sqrt{2})e^{i\theta/(\sqrt{2}v)}$.
\begin{equation}
|\partial\phi|^2 = \frac12(\partial R)^2 + \frac12(\partial\theta)^2 + \frac{\sqrt{2}R}{2v}(\partial\theta)^2 + \frac{R^2}{4v^2}(\partial\theta)^2 + \cdots
\end{equation}
Integrating out $R$ (heavy, $m_R^2 = 2\lambda v^2$): $s$-, $t$-, $u$-channel exchange $\;+\;$ contact term. Low-energy limit $s,t,u\ll m_R^2$:
\begin{equation}
i\mathcal{M} = \frac{i}{v^2}\Bigl[\frac{s^2}{s-m_R^2}+\frac{t^2}{t-m_R^2}+\frac{u^2}{u-m_R^2}\Bigr]
\;\xrightarrow{m_R\to\infty}\; -\frac{i}{v^2}(s+t+u + 3m_R^2) = -i(3\lambda),
\end{equation}
using $s+t+u=0$ (massless Goldstones).

\begin{equation}
\boxed{i\mathcal{M} = -i(3\lambda)\;\;\text{for both parametrizations}.}
\end{equation}

\end{document}
